\begin{abstract}
IEEE 802.11ax introduces Target Wake Time (TWT) and uplink OFDMA to reduce
station energy in dense Wi-Fi deployments.
The two mechanisms are usually configured separately, and no study measures
station energy when an access point controls both.
The difficulty is that a controller must adapt to changing traffic, yet online
learning is impractical in a production network, and a hand-tuned rule needs an
expert and is unlikely to be optimal.
Offline reinforcement learning avoids both: a policy is trained from logged
data alone, with no exploration on the live network.
We cast joint control as a 760-way action, selecting a TWT schedule, a
station-to-group assignment, and whether uplink OFDMA is enabled at each
decision, and train five offline algorithms on 100{,}235 logged transitions in
ns-3.44.
Implicit Q-Learning reaches $107.3$\,J against a tuned hand-designed
controller's $166.5$\,J, a $35.6\,\%$ reduction at higher throughput, and does
so consistently: the spread over ten training runs is $2.6$\,J. Four of the
five methods lose to that controller, and the energy of all five is ordered by
how often each enables uplink OFDMA.
\end{abstract}
